\begin{teilaufgaben}
\item
The function $u(x)$ is defined on the interval $\Omega = [0, \pi].$
The function satisfies in $\Omega$ 
\[
u''(x) - u(x) = 1
\]
and $u(0) = u(\pi) = 0$ on the boundary of $\Omega$.

Determine an approximation function $\tilde u(x)$ for $u(x)$ by
the method of Galerkin. 
Use the ansatz $\tilde u(x) = a_1 \cdot \sin(x) + a_2 \cdot \sin(3 \cdot x).$

\item
The function $u(x)$ is defined on the interval $\Omega = [0, 1].$
The function satisfies in $\Omega$ 
\[
- u''(x) = x
\]
and $u(0) = 1$ and $u(1) = 2$ on the boundary of $\Omega$.  

Determine an approximate value for $u(1/2)$ by the method of finite elements.

Use linear finite elements with $h = 1/2$.
\end{teilaufgaben}

\begin{loesung}
\begin{teilaufgaben}
\item
Plugging the ansatz into the differential equation
\[
u''(x) - u(x) - 1 = 0
\]
we obtain
\[
-2 \cdot a_1 \cdot \sin(x) - 10 \cdot a_2 \cdot \sin(3 \cdot x) - 1 \approx 0.
\]
Since
\[
\int_0^\pi \sin(x) \ dx = 2, \ \ \ \int_0^\pi \sin(3 \cdot x) \ dx = 2/3,
\]
as well as
\[
\int_0^\pi \sin(x)^2 \ dx = \pi/2, \ \ \ \int_0^\pi \sin(3 \cdot x)^2 \ dx = \pi/2, \ \ \  \int_0^\pi \sin(x) \cdot \sin(3 \cdot x) \ dx = 0,
\]
we obtain by Galerkin the linear system of equations
\begin{equation}
\left[\begin{array}{cc}  -\pi& 0 \\  0& -5 \cdot \pi  \end{array}\right]
\cdot
\left(\begin{array}{r} a_1 \\ a_2 \end{array}\right)
-
\left(\begin{array}{c} 2 \\ 2/3 \end{array}\right)
=
\left(\begin{array}{r} 0 \\ 0 \end{array}\right).
\tag{\bf{2P}}
\end{equation}
Therefore
\begin{equation}
(a_1, a_2) = -1/\pi \cdot (2, 2/15).
\tag{\bf{1P}}
\end{equation}

\item
The corresponding Ritz matrix is
\begin{equation}
R
=
\left[\begin{array}{rrr} 2 & -2 & 0 \\ -2 & 4 & -2  \\ 0 & -2 & 2 \end{array}\right],
\tag{\bf{1P}}
\end{equation}
whereas the Ritz vector is
\begin{equation}
\underline{r}
=
(*,1/4,*).
\tag{\bf{1P}}
\end{equation}
The vector of the knot variable is
\[
\underline{a} = (1, a, 2).
\]
Reduction of the Ritz system leads to  
\[
4 \cdot a - 6 = 1/4.
\]
Therefore we have
\begin{equation}
\tilde u(1/2) = 25/16.
\tag{\bf{1P}}
\end{equation}
\end{teilaufgaben}
\end{loesung}
